% Generated draft: requires mathematical and source review.
\documentclass[11pt]{article}
\usepackage[margin=25mm]{geometry}
\usepackage{fontspec}
\setmainfont{Latin Modern Roman}
\usepackage{amsmath,amssymb,mathtools,unicode-math}
\setmathfont{Latin Modern Math}
\usepackage{enumitem,needspace,xurl,hyperref}
\setlength{\parindent}{0pt}
\setlength{\parskip}{5pt}
\setlength{\emergencystretch}{3em}
\newcommand{\N}{\mathbb{N}}
\newcommand{\Z}{\mathbb{Z}}
\newcommand{\Q}{\mathbb{Q}}
\newcommand{\R}{\mathbb{R}}
\newcommand{\C}{\mathbb{C}}
\newcommand{\F}{\mathbb{F}}
\newcommand{\A}{\mathbb{A}}
\newcommand{\PP}{\mathbb{P}}
\newcommand{\E}{\mathbb{E}}
\newcommand{\eps}{\varepsilon}
\newcommand{\dd}{\,\mathrm{d}}
\begin{document}
\section*{UnsolvedMath ID: 20003298}
\textbf{Source problem code:} catalogue entry
\subsection{Definitions and mathematical statement}
This is a status paragraph for H. Oh's Question 24, not an independent question. Its opening joins the preceding record's phrase ``For the rank one groups,'' so the positive answer attributed to [26, 206, 155, 168, 104] has rank-one scope. The separate higher-rank claim concerns equidistribution of $\varepsilon$-separated closed geodesics from different homotopy classes [113]; it is not a statement about every compact flat. The final phrase ``one can also prove analog of the'' is incomplete.

The antecedent is recoverable from companion records. In a finite-volume locally symmetric setting $X=\Gamma\backslash G/K$, write $Y=\Gamma\backslash G$, and let $A$ be a Cartan subgroup. For each compact $A$-orbit $F\subset Y$, let $\mu_F$ be its Lebesgue orbit measure, $v_F=\mu_F(F)>0$, and $\bar\mu_F=\mu_F/v_F$. For the finite nonempty strict-cutoff family
\[
 \mathcal F_T=\{F:v_F<T\},
\]
Question 24 asks whether both
\[
 \nu_T=\frac{\sum_{F\in\mathcal F_T}v_F\bar\mu_F}
                  {\sum_{F\in\mathcal F_T}v_F},\qquad
 \bar\nu_T=\frac1{|\mathcal F_T|}\sum_{F\in\mathcal F_T}\bar\mu_F
\]
converge to normalized Haar measure on $Y$. The two normalizations are different.

The attached aid reports a sharp trimming lemma. For any finite nonempty family $I$ with positive weights $v_F$ and probabilities $\bar\mu_F$, let $\varnothing\ne J\subset I$ be a retained subfamily and form the corresponding volume-weighted and count-weighted probabilities $\nu_I,\nu_J$ and $\bar\nu_I,\bar\nu_J$. With total variation defined by $d_{\mathrm{TV}}(\rho,\sigma)=\sup_B|\rho(B)-\sigma(B)|$, it gives
\[
 d_{\mathrm{TV}}(\nu_I,\nu_J)
 \le\frac{\sum_{F\in I\setminus J}v_F}{\sum_{F\in I}v_F},\qquad
 d_{\mathrm{TV}}(\bar\nu_I,\bar\nu_J)
 \le\frac{|I\setminus J|}{|I|}.
\]
Both constants are sharp. Thus a selected-family volume-weighted equidistribution result transfers to the full family if the discarded volume fraction tends to zero; the corresponding count-weighted transfer is certified by a vanishing discarded count fraction. These are separate sufficient certificates and sharp worst-case bounds, not necessary conditions for convergence of every particular family. Nonvanishing discarded fractions can change limits, and separation or maximality alone does not establish either negligibility condition or equidistribution. The lemma therefore does not extend the source's restricted higher-rank result to the complete raw-volume-cutoff question.
\subsection{Short English statement}
This fragment reports a rank-one answer and a restricted higher-rank result for Question 24. Its supplied trimming bound shows how selected-family equidistribution transfers: discarded volume must be controlled for volume weighting, and discarded count for equal-orbit weighting.
\subsection{Sources}
\begin{itemize}
\item American Institute of Mathematics, \emph{Emerging applications of measure rigidity}, section 4.6, status paragraph following H. Oh's Question 24, \url{https://aimath.org/WWN/measrigid/measrigid.pdf}. Source references: [26, 206, 155, 168, 104] for rank one, and [113] for the restricted higher-rank statement. Full bibliography entries are not supplied.
\item Companion local records: permanent ID $20003295$, \texttt{AIM-OTHER-0102}, \texttt{notes[101]}, supplies the flat measures and both normalizations; ID $20003296$, \texttt{AIM-OTHER-0103}, \texttt{notes[102]}, states H. Oh's Question 24; ID $20003297$, \texttt{AIM-OTHER-0104}, \texttt{notes[103]}, supplies the opening rank-one qualifier. All are in canonical \texttt{aim-other-notes.json}.
\item Local UnsolvedMath registry, permanent ID $20003298$, source code \texttt{catalogue entry}, canonical \texttt{aim-other-notes.json}, \texttt{notes[104]}, dated 14 August 2026. Research attempt 1: context only, ``not a problem'' at run, proved rigor, low novelty confidence; machine-generated and not peer reviewed.
\end{itemize}
\end{document}
